\documentclass[UTF8,12pt]{ctexart}
\usepackage[a4paper,margin=2.4cm]{geometry}
\usepackage{amsmath,amssymb,booktabs,array,graphicx}
\usepackage{caption}
\captionsetup{font=small,labelfont=bf}
\setlength{\parindent}{2em}
\setlength{\parskip}{0.35em}

\title{问题一：基于季节画像、K-means 聚类与分层相关的蔬菜销量分析}
\author{}
\date{}

\begin{document}
\maketitle

\section{问题分析}

问题一要求同时说明各蔬菜品类及单品销量的分布规律与相互关系。销量序列具有非负、右偏和零值较多等特点，同时不同单品存在明显的季节上市规律。如果直接在全部日销量上计算一个相关系数，季节商品容易因大量零值被排除，且“共同客流造成的同涨”会与“销售结构中的此消彼长”混合。因此本文采用“零膨胀两部分布--月份与四季画像--K-means 季节型聚类--分层关系分析”的主线。旧版 MIC 与 Graphical Lasso 只作为附录中的严格条件关联对照，不承担本问主结论。

\section{数据与符号}

直接使用预处理后的日单品表和日品类表，时间范围为 2020 年 7 月 1 日至 2023 年 6 月 30 日。记单品 $i$ 在日期 $t$ 的销量为 $Y_{it}$（kg），题设品类 $c$ 的销量为
\begin{equation}
Y_{ct}=\sum_{i\in c}Y_{it}.
\end{equation}
销售年度定义为当年7月至次年6月，从而冬季12月至次年2月不会被年度边界拆开。单品须满足销售跨度不少于730天、正销售日不少于90天、累计销量不少于100 kg，不设置覆盖率门槛。由251个单品中筛得64个。

同一单品在首末销售记录之间的缺失日记为零销量；首日以前与末日以后视为不可观察。任意单品对的日相关只使用双方销售跨度的交集，避免将尚未上市或已经退市误记为零需求。

\section{模型建立}

\subsection{零质量与正销量条件分布}

对任一销量序列先估计零销量概率
\begin{equation}
\widehat p_0=\frac{1}{n}\sum_{t=1}^{n}\mathbf 1(Y_t=0).
\end{equation}
在 $Y_t>0$ 的样本上用极大似然法分别拟合 Normal、Lognormal、Gamma 和 Weibull 分布，以
\begin{equation}
\mathrm{AIC}=2K-2\ln L
\end{equation}
选择相对最优模型，并用 KS 检验评价绝对拟合质量。若AIC最小的参数分布仍在5\%水平被拒绝，则标记为KDE回退。该两部模型将“是否发生销售”与“发生销售后的销量大小”分开描述。

\subsection{月份与四季画像}

对每个实体补齐完整日历。记销售年度 $y$ 中月份 $m$ 的销量为 $Q_{iym}$，该年度总销量为 $Q_{iy}$，则月销量占比为
\begin{equation}
P_{iym}=\frac{Q_{iym}}{Q_{iy}},\qquad
P_{im}=\frac{1}{|\mathcal Y_i|}\sum_{y\in\mathcal Y_i}P_{iym},
\end{equation}
其中 $\mathcal Y_i$ 只包含年度总销量大于0的销售年度，因此 $\sum_{m=1}^{12}P_{im}=1$。月活跃率定义为
\begin{equation}
A_{iym}=\frac{\text{月份 }m\text{ 中正销量日数}}{\text{月份 }m\text{ 的日历日数}},\qquad
A_{im}=\frac{1}{3}\sum_y A_{iym}.
\end{equation}
月季节指数为月均日销量与全年月均日销量平均值之比。四季按春季3--5月、夏季6--8月、秋季9--11月、冬季12--2月汇总。

\subsection{季节型 K-means 聚类}

每个单品采用24维特征
\begin{equation}
\mathbf x_i=(P_{i1},\ldots,P_{i12},A_{i1},\ldots,A_{i12})^{\mathsf T},
\end{equation}
各维标准化后执行K-means：
\begin{equation}
\min_{C_1,\ldots,C_k}\sum_{r=1}^{k}\sum_{i\in C_r}
\left\|\mathbf z_i-\boldsymbol\mu_r\right\|_2^2.
\end{equation}
比较 $k=2,\ldots,8$ 的轮廓系数、Calinski--Harabasz 指数、Davies--Bouldin 指数、最小簇规模和100次完整销售年度块重采样ARI。综合可解释性、无微小簇和稳定性选择 $k=5$，再依据簇中心的峰值月份与四季占比一对一命名为春季型、初夏型、夏秋型、冬季型和常年型。

\subsection{四类关系指标与可信度}

为避免用一个系数概括不同含义，同时计算：
\begin{enumerate}
  \item 12个月季节指数的Pearson相关，用于判断旺淡季曲线是否一致；
  \item 日销量取 $\log(1+Y)$ 并减去年月均值后的分季节Spearman相关，用于判断同一季节内是否同向波动；
  \item 每日销量份额的分季节Spearman相关，用于区分总量同涨与销售结构竞争；
  \item 正销售日集合的Jaccard系数
  \begin{equation}
  J_{ij}=\frac{|A_i\cap A_j|}{|A_i\cup A_j|},
  \end{equation}
  用于衡量两者是否经常同期销售。
\end{enumerate}

六品类15对、五季节簇10对和64单品2016对均完整保留。对品类与簇关系完整执行200次7日移动块Bootstrap；单品层对达到 $|\rho|\ge0.30$ 的论文候选关系执行相同Bootstrap，其余关系保留点估计但不伪造置信区间。所有$p$值按指标和季节分别作BH校正。正文将 $|\rho|\ge0.30$ 且95\%区间不跨0的关系称为“较明确关系”，但不据此删除弱关系。

\section{结果与解释}

\subsection{销量分布和季节规律}

六品类的零销量比例均较低，正销量分布全部通过KS检验。其中水生根茎类、花叶类、花菜类和茄类以Gamma分布的AIC最小，辣椒类和食用菌以Lognormal分布最优。64个单品与6个品类共70个对象中，45个接受参数分布，25个回退KDE，说明单品层的分布异质性明显高于品类层。

品类峰值月份显示：水生根茎类和食用菌在1月达到峰值，辣椒类在2月达到峰值；花叶类和花菜类在8月达到峰值，茄类在7月达到峰值。水生根茎类的季节集中度最高，冬季销量占比明显高于夏季；花叶类和辣椒类的全年分布相对均衡。

\subsection{聚类结果}

修正后的 $k=5$ 轮廓系数为0.23，CH指数为15.79，DB指数为1.48，完整销售年度块重采样平均ARI为0.57。五簇规模分别为7、8、11、17和21，没有单例或极小簇。按语义排序，初夏型11个、冬季型7个、常年型21个、春季型8个、夏秋型17个；其主季销量占比分别约为62.71\%、70.81\%、33.19\%（夏季为常年型的相对最高季）、54.77\%和41.15\%。这些结果说明64个长期销售单品仍存在清晰的季节供给差异，但ARI处于中等水平，簇边界不是绝对分类。

\subsection{品类和簇间关系}

六品类全年去年月效应后的销量相关均表现为同向，其中花叶类--辣椒类为0.69，花叶类--食用菌为0.69，辣椒类--食用菌为0.65，Bootstrap区间均不跨0。这说明在共同客流、节假日或短期经营节奏变化下，多数品类总销量会同时增减。

但销量份额给出不同信息。例如花叶类--食用菌全年份额相关为$-0.54$，水生根茎类--花叶类为$-0.46$。因此“总量同涨”并不等于结构互补：两类商品可以在客流上升时同时卖得更多，同时在当日蔬菜总销量构成中此消彼长。份额相关受到各份额和为1的组成约束，只解释为销售结构竞争，不解释为价格替代或因果替代。

簇间关系进一步表明，常年型与冬季型全年销量相关为0.39，常年型与夏秋型为0.32，初夏型与常年型为0.31，均具有稳定同向波动。常年型--夏秋型全年份额相关为$-0.72$，初夏型--冬季型为$-0.79$，说明不同季节供给组在总销量结构上具有明显轮换。

\subsection{单品层关系}

完整机器表保留2016对单品关系。全年销量相关中有18对满足 $|\rho|\ge0.30$ 且区间不跨0。较强的正向关系包括“西兰花--净藕(1)”（$\rho=0.44$，95\%区间$[0.38,0.52]$）、“西兰花--芜湖青椒(1)”（0.44，$[0.38,0.51]$）和“七彩椒(1)--红灯笼椒(1)”（0.43，$[0.33,0.52]$）。较明确的负向关系为“西峡花菇(1)--西峡香菇(1)”（$-0.33$，$[-0.43,-0.23]$）。这些关系表示共同销售节奏或结构分化，并不能直接推导购物篮互补、价格弹性或因果替代。

\section{模型评价}

模型的优点是保留季节性单品，以月份画像直接刻画题意；k的选择同时考虑分离度、簇规模和跨年度稳定性；总量相关、份额相关和活跃期重合分别回答不同的“关系”；完整2016对母表保证弱关系和跨品类关系不被遗漏。主要局限是无销售记录仍可能混合无需求、缺货和未陈列；三个销售年度对长期稳定性判断仍有限；份额相关存在组成数据的天然负相关约束；所有关系均为统计关联而非因果效应。

综上，问题一的最终解释不是“蔬菜之间基本没有关系”，而是：蔬菜销量具有显著季节分型；多数品类在短期总量上共同涨落，但在销量结构中可能竞争；单品关系具有明显异质性，应在品类、季节簇和具体单品三个尺度分别解释。

\end{document}
